\uuid{3tj7}
\titre{Nombres premiers}
\auteur{bodin ; croizat ; sacre}
\organisation{exo7}
\licence{CC BY-NC-SA 4.0}
\datecreate{2022-03-08}
\niveau{L1}
\module{Algèbre}
\chapitre{Arithmétique dans Z}
\sousChapitre{Nombres premiers, nombres premiers entre eux}
\difficulte{3}
\theme{Nombres premiers}
\exo7qcmid{hilisit.arithmetique.037}

\contenu{
\begin{question}
\qclassification{103.04}
    Soient $a = 5^3 \times 11^2 \times 13^5 \times 19$
    et $b = 5^5 \times 7^4 \times 11 \times 19$
    Quelles sont les affirmations vraies ?
    \begin{answers}
        \bad{$\pgcd(a,b) = 5^3 \times 7^4 \times 11 \times 13^5 \times 19$}
        \bad{$\pgcd(a,b) = 5 \times 11 \times 19$}
        \good{$\ppcm(a,b) = 5^5 \times 7^4 \times 11^2 \times 13^5 \times 19$}
        \bad{$\ppcm(a,b) = 5^5 \times 11^2 \times 19$}
    \end{answers}
    \begin{explanations}
    Pour le pgcd, on garde le plus petit exposant des décompositions de $a$ et $b$ ; pour le ppcm, on garde le plus grand exposant. \\
     $\pgcd(a,b) = 5^3 \times 11 \times 19$ \\
     $\ppcm(a,b) = 5^5 \times 7^4 \times 11^2 \times 13^5 \times 19$
    \end{explanations}
\end{question}
}
